% Propietario conservado: /Users/ruben/Documents/New project/output/EXTREMA_MEDIA_RAZON_RECIPROCIDAD_ES_EN_20260918/ARTICULO/ES/source/nuclear/04.tex
\chapter{Exceptional incidence and transport of forms}
\label{x_reciprocidad:nuclear:4}
\section{Generation of the record and conservation of incidence}

\subsection{Dodecaphasic record of the enriched history}

Section~\ref{x_reciprocidad:subsec:k-registro-integral} defines

\[
\mathcal R_{12}(x)=((E_m,\tau_m,\sigma_m,\kappa_m,\Lambda_m))_{m=1}^{12}.
\]

It preserves the window, subpath, orientation, carry, and incidence ledger. Section~\ref{x_reciprocidad:subsec:k-terminal} fixes

\[
x_{\rm term}=(\mathcal U_{108}\cdots\mathcal U_1)x_{\rm seed},
\quad \mathcal U_t=\mathsf{Upd}_t\mathsf{Tra}_t\mathsf{Sel}_t,
\]

\[
R_{\rm sgn}=\Pi_H\mathcal E_{108}^{90,120}\mathcal R_{12}:
\mathcal X_{\rm TPK}^{\rm term,mem}\to\mathbb Z^{12}.
\]

The following algebraic reconstruction begins with this typed output. It does not replace evaluation of the preceding events or obtain their coordinates by imposing \(K\) or alpha as a target.

\subsection{Integral normalization and isometric normalizations}

For the step-three orbits \((1,4,7,10)\), \((2,5,8,11)\), \((3,6,9,12)\), Lemma~\ref{x_reciprocidad:lem:k-hadamard} fixes the symmetric matrix \(H=H_{12}\) and proves \(H^2=4I\).

\[
\mathcal L_H=\{u\in\mathbb Z^{12}:Hu\equiv0\pmod4\},
\qquad \mathscr K(u)=\tfrac14Hu,
\qquad \mathscr K^{-1}(k)=Hk.
\]

This is an integral isomorphism between the sublattice of signed records and \(\mathbb Z^{12}\). The factor \(1/4\) implements inversion, not a Euclidean isometry.

As an immediate consequence of the same identity, \(J_H=H/2\) satisfies \(J_H^*J_H=I\): this is the real isometric transport. If \(u=Hk\), then \(u/2=J_Hk\) and \(\|u\|^2=4\|k\|^2\). Confusing \(H/4\) with \(H/2\) would change the balance by a factor of four.

In \cref{x_reciprocidad:exc:entrelazador,x_reciprocidad:exc:isometria}, the constructed design determines

\[
I:\mathbb R^{12}\to\mathbb R^{132},\qquad
(Iv)(H')=\sum_{p\in H'}v_p,\qquad I^*I=36I_{12}+30J_{12}.
\]

The coefficients arise from the incidences: each point belongs to 66 hexads, and each pair to 30. On \(V=\mathbf1^\perp\), the term \(J_{12}\) vanishes, so \(U_W=I/6:V\to E_{\rm inc}\) is an isometry. The factor \(1/6\) normalizes this screen; it is not the inverse map for \(K\).

\subsection{Prefix refinement and event partitions}

Section~\ref{x_reciprocidad:exc:doble-lectura} constructs causal frames \(f_{j+1}=g_j(x)f_j\) and prefix encodings \(Q_n\). It proves

\[
r_{n,m}Q_n=Q_mp_{n,m},\qquad
\gamma_{A^f}(wf^{-1})=\gamma_A(w)(f^{-1}\oplus f^{-1}).
\]

The proof uses the fact that a prolongation preserves preceding words, transports, and frames. The inverse limit produces \(Q_\infty\). Recovery of the visible word follows from the first component of \(\gamma\); the text expressly distinguishes this scope from recovery of all deep memory and passage to supports.

The additive and affine composition of the record, described in \cref{x_reciprocidad:subsec:k-registro-integral,x_reciprocidad:iv:cont:concatenacion}, proves a second conservation property: with compatible prospective contributions \((\delta w_t,\delta q_t)\), record inversion is additive. Concatenating segments or subdividing the same list of events gives the same final record. The version with nontrivial transport uses the affine cocycle. Preserving the total under a change of partition does not assert that the record stops growing when new events are added.

\section{K–Witt transport of norms, operators, and balances}\label{x_reciprocidad:nuc04:kwitt}

A formalization of the recovered results is assembled here, together with the algebraic proof below. This assembly attributes no new priority to its starting identities.

Let \(P_0=I-J_{12}/12\), \(V=P_0\mathbb R^{12}\), \(W_H=J_HV\), and

\[
B:=U_WJ_H|_{W_H}:W_H\longrightarrow E_{\rm inc}.
\]

Since \(J_H^2=I\), this is an isometric isomorphism between eleven-dimensional spaces. For a record \(k\), its normalized signed coordinate is \(J_Hk\). If \(\widehat P_0=J_HP_0J_H\),

\[
B\widehat P_0(J_Hk)=U_WP_0k,
\]

\[
\|\widehat P_0(J_Hk)\|^2
=\|U_WP_0k\|^2
=\sum_{p=1}^{12}k_p^2-\frac{(\sum_pk_p)^2}{12}.
\]

No constant value has been introduced to obtain this equality: it follows from \(H^2\) and the incidence Gramian. The removed mean must be retained separately to recover the whole of \(k\).

For \(g\in\operatorname{Aut}(\mathcal H)\), let \(\rho_{12}(g)\) be its point permutation and \(\rho_{\rm hex}(g)\) its hexad permutation. The action on signed coordinates is

\[
\widehat\rho_H(g)=J_H\rho_{12}(g)J_H
=\tfrac14H\rho_{12}(g)H.
\]

It is orthogonal, preserves \(W_H\), and also preserves the sublattice \(\mathcal L_H\): if \(u=Hk\), its image is \(H\rho_{12}(g)k\). The explicit intertwining in \cref{x_reciprocidad:exc:entrelazador} gives

\[
B\widehat\rho_H(g)=\rho_{\rm hex}(g)B.
\]

\begin{proof}
Substitute the definitions, cancel \(J_H^2\), and apply \(U_W\rho_{12}(g)=\rho_{\rm hex}(g)U_W\). The integral record, centred inner product, and incidence action are thus preserved, each on its own domain.
\end{proof}

For an operator \(A\) already constructed on \(W_H\), one obtains \(A_{\rm inc}=BAB^*\). By substitution,

\[
BA=A_{\rm inc}B,\quad
\langle z,Az\rangle=\langle Bz,A_{\rm inc}Bz\rangle.
\]

Therefore, if a transport \(C\) on \(W_H\) preserves \(A\), \(C^*AC=A\), its transport \(C_{\rm inc}=BCB^*\) preserves \(A_{\rm inc}\). The condition that a particular \(C\) be a publication of TPK dynamics remains in the causal composition producing it. Isometry does not by itself transform an incidence permutation into the nonadic cycle.

Commutation need not be imposed to compare transported charts. If \(R\) is unitary, \(C'=RCR^*\) and \(T(C)=(8I+C)/9\), then \(T(C')R=RT(C)\) by substitution. Together with \(A'=RAR^*\), the balance in one chart is conjugated to that in the other. Commutation \([R,C]=0\) corresponds only to the additional problem of keeping that same seam \(C\) fixed. This distinction agrees with the lattice naturality of Section~\ref{x_reciprocidad:nuc04:c3}.

One available concrete finite action is the star selected by the flag:

\[
\sigma_{B_K}=(1\ 3)(2\ 11)(5\ 7)(8\ 12).
\]

Section~\ref{x_reciprocidad:exc:estrella} proves its construction and distinguishes it from the group generated by the 495 stars. The individual star gives \(C_2\); the family gives \(M_{12}\). Transporting the marking and origin gives covariance between frames. If a symmetry of a fixed frame is sought, the appropriate domain is its stabilizer, not automatically all of \(M_{12}\). If it must additionally commute with \(C\), the subgroup satisfying the specific equation \([\widehat\rho_H(g),C]=0\) is used.

\section{Order-three action on the lattice fibres}\label{x_reciprocidad:nuc04:c3}

The construction in \cref{x_reciprocidad:nuclear:fibras-app-witt,x_reciprocidad:nuclear:alineamiento-witt} produces twelve \(A_2\) fibres labelled by \((i,\mu,\varepsilon)\in\mathbb Z/3\mathbb Z\times\{0,1\}\times\{0,1\}\). The alignment \(\lambda(i,\mu,\varepsilon)=4i+2\mu+\varepsilon\pmod{12}\) is bijective. If \(C_{A_2}\) is the Coxeter operator and \(R\) the reflection satisfying \(RC_{A_2}R=C_{A_2}^{-1}\), the trivializations and transported frames fix

\[
g_b=\iota_b^{-1}P_bC_{A_2}P_b^{-1}\iota_b,
\quad T_{i\mu\varepsilon}=\iota_{\lambda(i,\mu,\varepsilon)}^{-1}P_{\lambda(i,\mu,\varepsilon)}R^\mu,
\]

\[
g_\lambda T_{i\mu\varepsilon}
=T_{i\mu\varepsilon}C_{A_2}^{(-1)^\mu}.
\]

The direct sum \(T_\lambda:W_{\rm APP}\otimes\mathbb C\to\bigoplus_{b=1}^{12}(A_{2,b}\otimes\mathbb C)\) is a \(C_3\)-equivariant isomorphism. The proof consists of conjugation and the bijection of the twelve labels: it preserves pair, sheet, and orbit. It is not a coincidence of dimensions.

The constructions in \cref{x_reciprocidad:km:palabra-total,x_reciprocidad:km:estabilidad-vecino} use \(q^\star=(1,1,1,1,1,1)\) and its full-support word \(c^\star=(q^\star,q^\star A_W)\) to orient the twelve blocks. The resulting \(g_c\) has order three, has no nonzero fixed vectors, and preserves the Niemeier gluing. For the incidence vector \(v\) of squared norm \(54\),

\[
N_0=\ker(\langle\cdot,v\rangle\bmod3),\qquad
N_v=N_0+\mathbb Z\frac v3,
\]

\[
y_*=(g_cv-v)/3\in N_0,\qquad
z_*=(g_c^{-1}v-v)/3\in N_0.
\]

The proof identifies their discriminant classes as \(c^\star\) and \(-c^\star\), and their pairings with \(v\) as \(-27\). Hence \(g_cN_0=N_0\) and \(g_c(v/3)=v/3+y_*\), giving \(g_cN_v=N_v\). The preserved structure includes the gluing, congruence kernel, and class generating the Leech neighbour.

Prolongation to the 24 full-support orientations and the lattice lift of a monomial code automorphism \(h\) are developed in Section~\ref{x_reciprocidad:paper:24-orientaciones}. With all marks transported,

\[
g_{hu}=\widetilde h\,g_u\,\widetilde h^{-1}.
\]

This \(C_3\) is an admissible group genuinely intertwined with the APP fibres and exceptional structure. It is not identified here with \(C_9\) or its quotient by a mere relation of orders. Nor is the lattice lift replaced by the ternary matrix without its action on the \(A_2\) fibres.

\section{Refinement of energy and its differential}
\label{x_reciprocidad:paper:refinamiento-energia}

The energy operator of \eqref{x_reciprocidad:vii:eq:energia-memoria-discreta} has the form \(D_{9,m}^*W_mD_{9,m}\), with \(U_a\) transporting orientation and carry. The domains in \cref{x_reciprocidad:vii:subsec:funcional-memoria} are

\[
\mathcal H_m=\bigoplus_vH_v,\quad
\mathcal K_m=\bigoplus_{a\in\mathcal E_m^+}H_{t(a)},\quad
(D_mf)_a=f_{t(a)}-U_af_{s(a)},\quad
E_m=\langle D_mf,\mathsf W_mD_mf\rangle.
\]

The fixed maps \(J_m\) and \(K_m\) satisfy, for the family \(C^1\) considered,

\[
D_{m+1}(t)J_m=K_mD_m(t),\qquad
K_m^*\mathsf W_{m+1}(t)K_m=\mathsf W_m(t).
\]

The result in \cref{x_reciprocidad:vii:thm:refinamiento-respuesta} proves \(E_{m+1}(J_mf)=E_m(f)\) by substitution, and equality of differentials follows by differentiating both identities. For parameter-dependent refinement, the term \(D_{m+1}\dot J_mf\) is also retained.

With \(v_a=U_af_s\) and \(q=\mathsf W_mD_mf\), the response to \(\dot U_a=A_aU_a\) is

\[
\mathfrak j_m(A)=-2\operatorname{Re}\sum_a\langle q_a,A_av_a\rangle,
\quad \mathcal J_a=v_aq_a^*-q_av_a^*.
\]

The transport product \(U_\gamma=U_N\cdots U_1\) satisfies

\[
\dot U_\gamma U_\gamma^{-1}
=\sum_{k=1}^N B_kA_kB_k^{-1},\qquad B_k=U_N\cdots U_{k+1}.
\]

This formula retains the order of each incidence and transports its covector contragrediently. Corollary~\ref{x_reciprocidad:vii:cor:refinamiento-corriente-conexion} composes this differential with connection variations and proves, for compatible refinements of the action,

\[
s_m=P_m^{\mathsf T}s_{m+1},\qquad
M_m\sigma_m=P_m^{\mathsf T}M_{m+1}\sigma_{m+1}.
\]

Here, \(P_m\) carries connection variations, and \(M_m\) is the cellular pairing; they are neither the mean-zero projector nor the Hadamard matrix.

\subsection{Assembly with incidence}

In a realization of fibres \(W_H\) whose frame actions are \(\widehat\rho_H(g_a)\), conjugation by \(B\) carries every transport, field, and weight to \(E_{\rm inc}\). By the intertwining proved in Section~\ref{x_reciprocidad:nuc04:kwitt}, the difference operators satisfy \(D_{\rm inc}B_{\rm vertices}=B_{\rm edges}D_H\). Transporting \(J_m\), \(K_m\), and \(\mathsf W_m\) as well, the two preceding refinement squares remain the same conjugated squares. By substitution, they preserve energy; for fixed \(B\), they also preserve the differential and covectors.

The balance law and its covectors are transported to the exceptional representation while retaining compatibility with refinement. Each map uses the transports of the specified dynamics and the subspaces they preserve. The variation domain is defined in \cref{x_reciprocidad:vii:subsec:funcional-memoria}.
